\chapter{Division}\label{ch:g4:division}

Grade 3 divided with the tables and small \omterm{def:g3:division:remainder}{remainders}
(\cref{ch:g3:division}); this chapter installs the full written
method --- long division --- which handles any number, digit group by
digit group.

\section{The long division method}

\begin{method}[Long division by a one-digit number]\label{met:g4:division:method}
To divide $749$ by $6$:
\begin{enumerate}
  \item take the leftmost digit, $7$ (hundreds): $6$ goes $1$ time
        into $7$; write $1$ in the \omterm{def:g3:division:remainder}{quotient}, subtract
        $6$, \omterm{def:g3:division:remainder}{remainder} $1$;
  \item \emph{bring down} the next digit, $4$: now divide $14$
        (tens): $6 \times 2 = 12$; write $2$, \omterm{def:g3:division:remainder}{remainder} $2$;
  \item bring down the $9$: divide $29$: $6 \times 4 = 24$; write
        $4$, \omterm{def:g3:division:remainder}{remainder} $5$;
  \item no digit left: the \omterm{def:g3:division:remainder}{quotient} is $124$, the \omterm{def:g3:division:remainder}{remainder} $5$.
\end{enumerate}
Check: $6 \times 124 + 5 = 744 + 5 = 749$, and $5 < 6$.
\end{method}

\begin{omfigure}
\begin{tikzpicture}[scale=0.95]
  % digit columns at x = 0.2, 0.8, 1.4; every digit under its place
  \node[font=\normalsize] at (0.2,3.0) {$7$};
  \node[font=\normalsize] at (0.8,3.0) {$4$};
  \node[font=\normalsize] at (1.4,3.0) {$9$};
  % divisor bracket and quotient
  \draw[thick] (1.8,2.0) -- (1.8,3.35);
  \draw[thick] (1.8,2.62) -- (3.7,2.62);
  \node[font=\normalsize] at (2.35,3.0) {$6$};
  \node[font=\normalsize, omDef] at (2.35,2.2) {$1$};
  \node[font=\normalsize, omDef] at (2.95,2.2) {$2$};
  \node[font=\normalsize, omDef] at (3.55,2.2) {$4$};
  % step 1: 7 - 6 = 1
  \node[font=\normalsize, omThm] at (-0.35,2.4) {$-$};
  \node[font=\normalsize, omThm] at (0.2,2.4) {$6$};
  \draw (-0.55,2.1) -- (0.55,2.1);
  \node[font=\normalsize] at (0.2,1.8) {$1$};
  \node[font=\normalsize] at (0.8,1.8) {$4$};
  \draw[-Stealth, omExo, thin] (0.8,2.75) -- (0.8,2.05);
  % step 2: 14 - 12 = 2
  \node[font=\normalsize, omThm] at (-0.35,1.2) {$-$};
  \node[font=\normalsize, omThm] at (0.2,1.2) {$1$};
  \node[font=\normalsize, omThm] at (0.8,1.2) {$2$};
  \draw (-0.55,0.9) -- (1.15,0.9);
  \node[font=\normalsize] at (0.8,0.6) {$2$};
  \node[font=\normalsize] at (1.4,0.6) {$9$};
  \draw[-Stealth, omExo, thin] (1.4,2.75) -- (1.4,0.85);
  % step 3: 29 - 24 = 5
  \node[font=\normalsize, omThm] at (0.25,0.0) {$-$};
  \node[font=\normalsize, omThm] at (0.8,0.0) {$2$};
  \node[font=\normalsize, omThm] at (1.4,0.0) {$4$};
  \draw (0.05,-0.3) -- (1.75,-0.3);
  \node[font=\normalsize] at (1.4,-0.6) {$5$};
  % commentary
  \node[right, font=\small, omExo, align=left] at (4.5,1.4)
    {$7 - 6 = 1$, bring down the $4$;\\
     $14 - 12 = 2$, bring down the $9$;\\
     $29 - 24 = 5$: stop.\\[3pt]
     quotient $124$, remainder $5$\\
     check: $6 \times 124 + 5 = 749$};
\end{tikzpicture}

{\small The long division of $749$ by $6$, written in full: each
subtraction appears, every digit stays in its column, and the gray
arrows show the digits being brought down.}
\end{omfigure}

\begin{example}[A zero in the quotient]\label{ex:g4:division:zero}
Divide $618$ by $3$: hundreds, $6 \div 3 = 2$; tens, bring down the
$1$: $3$ goes $0$ times into $1$ --- \emph{write the} $0$! ---
\omterm{def:g3:division:remainder}{remainder} $1$; units, bring down the $8$: $18 \div 3 = 6$. \omterm{def:g3:division:remainder}{Quotient}
$206$, \omterm{def:g3:division:remainder}{remainder} $0$. Forgetting the middle \omterm{def:g1:counting:zero}{zero} (writing $26$)
is the classic mistake; the check
$3 \times 26 = 78 \neq 618$ catches it immediately.
\end{example}

\begin{example}[Estimating the quotient's size]\label{ex:g4:division:size}
Before dividing $749$ by $6$, frame the answer:
$6 \times 100 = 600$ and $6 \times 200 = 1\,200$, so the \omterm{def:g3:division:remainder}{quotient} is
between $100$ and $200$ --- it will have three digits. This one-line
estimate prevents most misplaced-digit errors.
\end{example}

\section{Choosing what the question asks}

\begin{method}[Quotient, remainder, or quotient plus one]\label{met:g4:division:interpret}
As in \cref{met:g3:division:interpret}, the story decides:
\begin{enumerate}
  \item ``how many full \dots'' or ``how many each'': the \omterm{def:g3:division:remainder}{quotient};
  \item ``how many left'': the \omterm{def:g3:division:remainder}{remainder};
  \item ``how many containers needed for all'': the \omterm{def:g3:division:remainder}{quotient}, plus
        one if the \omterm{def:g3:division:remainder}{remainder} is not \omterm{def:g1:counting:zero}{zero}.
\end{enumerate}
\end{method}

\begin{example}\label{ex:g4:division:problem}
$530$ books must be packed in boxes of $8$.
\[
530 = 8 \times 66 + 2 .
\]
How many \emph{full} boxes? $66$. How many books are not in a full
box? $2$. How many boxes to pack \emph{all} the books? $67$.
Three questions, one division.
\end{example}

\begin{example}[Fair sharing of money]\label{ex:g4:division:money}
Four friends share the cost of a $92$ gift equally:
$92 \div 4 = 23$ exactly (\omterm{def:g3:division:remainder}{remainder} $0$). Each pays $23$. When the
\omterm{def:g3:division:remainder}{remainder} is \omterm{def:g1:counting:zero}{zero}, the division ``comes out even'' and the sharing
is perfectly fair.
\end{example}

\section{Exercises}

\begin{exercise}[$\star$]\label{exo:g4:division:1}
Frame each \omterm{def:g3:division:remainder}{quotient} as in \cref{ex:g4:division:size} (between which
two ``\omterm{def:g4:numbers:round}{round}'' numbers?), then compute the long division:
$96 \div 4$; \quad $87 \div 5$.
\end{exercise}

\begin{exercise}[$\star$]\label{exo:g4:division:2}
Compute the long divisions and check each one:
$672 \div 4$; \quad $925 \div 7$; \quad $804 \div 6$.
\end{exercise}

\begin{exercise}[$\star$]\label{exo:g4:division:3}
Careful with the \omterm{def:g1:counting:zero}{zero} (see \cref{ex:g4:division:zero}):
$816 \div 4$; \quad $420 \div 7$; \quad $2\,512 \div 5$.
\end{exercise}

\begin{exercise}[$\star$]\label{exo:g4:division:4}
Write the check equality ($a = b \times q + r$) for: $85 \div 9$;
\quad $1\,000 \div 3$.
\end{exercise}

\begin{exercise}[$\star$]\label{exo:g4:division:5}
$156$ eggs are put in boxes of $6$. How many boxes are filled?
\end{exercise}

\begin{exercise}[$\star$]\label{exo:g4:division:6}
A ribbon of $250$ cm is cut into pieces of $8$ cm. How many pieces,
and what length is left over?
\end{exercise}

\begin{exercise}[$\star$]\label{exo:g4:division:7}
$375$ students go to a show in buses of $52$ seats. How many buses
are needed? (Which case of
\cref{met:g4:division:interpret} is this? Divide by $52$ using
multiples: $52 \times 7 = 364$.)
\end{exercise}

\begin{exercise}[$\star$]\label{exo:g4:division:8}
Three friends share $234$ marbles equally. How many marbles each?
Check with a multiplication.
\end{exercise}

\begin{exercise}[$\star$]\label{exo:g4:division:9}
Find the dividend: the division by $7$ gives \omterm{def:g3:division:remainder}{quotient} $58$ and
\omterm{def:g3:division:remainder}{remainder} $3$.
\end{exercise}

\begin{exercise}[$\star\star$]\label{exo:g4:division:10}
A librarian shelves $438$ books, $9$ per shelf. Shelves come in
bookcases of $6$ shelves. How many shelves are needed? How many
bookcases must be bought? (Two divisions, both with a
``plus one'' question.)
\end{exercise}

\begin{exercise}[$\star\star$]\label{exo:g4:division:11}
In a division by $8$, the \omterm{def:g3:division:remainder}{quotient} equals the \omterm{def:g3:division:remainder}{remainder}. What are
the possible dividends? List them all. (The \omterm{def:g3:division:remainder}{remainder} must stay
below $8$.)
\end{exercise}
